\section{Weak-C2 R3 with a positive-unit $1-r$ is admitted}
\label{sec:r18-r3-eunit}

\begin{theorem}
Keep the entire original physical stationary weak-C2 source, $d=1$,
$\lambda_{\rm event}=\beta=0$, every strict gap and nonnegative
coefficient, and every original residual intersection. Suppose
$B,z,y,l=1-h$ are positive infinitesimals, $t=y/B$ a positive bounded
unit, and $e=1-r$ a positive bounded unit. Then the full original
parent is a finite tensor SOS for every independent projective testing
pair and every real or complex state. These are precisely the
positive-unit-$e$ sources of R3. The remaining $e$-infinitesimal
part of R3 is not asserted here.
\end{theorem}

\paragraph{Source estimates including all simultaneous openings.}
Put $a=z/B$, $g=1-ae+haB$ and $j=1+aBr-l/t$. The positivity
of $g$ and the unit lower bound for $e$ imply that $a$ is bounded.
Moreover $h\sim1$ and $j\sim1$. In the following formulas variables
on the right denote their real residues, $e=1-r>0$ and
$g=1-ae\ge0$. Direct joint regular reduction of the unchanged WE1
formulas gives
\[
 (nr)_0={e\over2(1+a)},\qquad
 (mg)_0={a^2e\over2(1+a)},\qquad
 A_0={t(1+ar)\over2(1+a)},\qquad
 D_0={1+ar\over2t(1+a)},                         \tag{R3E1}
\]
\[
 (Cl)_0={(1+ar)(1+t)\over2(1+a)},\qquad
 \alpha_{*,0}:=(\alpha/B^2)_0
       ={(1+ar)^2(1+t)\over2(1+a)}.              \tag{R3E2}
\]
Thus $A,D,Cl,\alpha_*$ are positive bounded units, $C$ is unbounded,
and $n$ has a positive real lower bound, even when $r$ is
infinitesimal and $n$ is unbounded. In particular literal $n=0$
is impossible on this source domain, by the positive first formula
in R3E1; no coefficient boundary was discarded. Write
$N=n+1/2$, so $N^{-1}$ is bounded. If $g$ is a unit, then
$\mu=m/a^2$ is a positive unit with
\[
              \mu_0={e\over2g(1+a)}.             \tag{R3E3}
\]
If $g$ is infinitesimal, $a$ is a unit, $m$ is unbounded and $mg$
is a positive unit. These joint statements use only the displayed
products: the WE1 denominators after clearing $r,g,l$ have unit
factors $t,h,1-hB,\Delta$ with $\Delta_0=1+a$.
There is no comparison among the $B,l,r,g$ openings. The original
positive-infinitesimal total-scale condition follows already from
$C$ unbounded.

\paragraph{The full local matrix with its divergent C pivot retained.}
Use the exact complete WA3--WA8 parent identity in the roots
\[
 \xi=X+ruZ,\quad \eta=Y-huW,\quad
 R_1=uW-vZ/t,\quad R_2'=avW-uZ,\quad T=uvW-BZ.
\]
Normalize its first two rows by
$\widehat\xi=\sqrt N\xi$ and $\widehat\eta=\sqrt C\eta$.
First suppose $g$ is a unit. The complete resulting five-root
matrix has bounded pure-source polynomial coefficients in $u,v$.
Here is the coefficient check, including the potentially singular
source openings. The hard entries use $n/N$, $C/C$, and the bounded
coefficients $A,D,m,\alpha$, together with bounded $N^{-1/2}$ and
infinitesimal $C^{-1/2}$. All scalar occurrences of $n$ are through
$nr^2$, and all first-normal occurrences are through $nr$, both
bounded. All non-hard occurrences of $C$ are through $Cl$ or
$Cl^2$, both bounded. Thus the entire $\eta$ normal row becomes
infinitesimal after division by $\sqrt C$. This includes the full
WA7 terms, not merely their zero-angle values.

The WE5 scaled-$R_2'$ formulas bound every coefficient through
$a\to0$ except the apparently singular
\[
 Q_*={b_2-\alpha_*\over a},\qquad
 b_2=2a^2\mu h+Cl^2+Ah^2+(B^2\alpha_*+1/2)h^2.  \tag{R3E4}
\]
The inherited exact WE3 numerator vanishes at $a=0$ for arbitrary
$B,h,r,t$. Its quotient has a unit denominator here too: the $l$
denominator cancels in $Cl^2$ before division, and the remaining
factors $g,h,t,j,1-hB,\Delta$ are units on this $g$-unit chart.
Consequently $Q_*$ is bounded jointly through $a,l,r\to0$.
This bounds the entire scaled scalar block and both normal rows,
without any $a/l$ or $l/a$ hypothesis.

At $u=v=0$ the full residue has separate pivots $1,A_0,\alpha_{*,0}$
on $\widehat\eta,R_1,T$. Its $\widehat\xi,R_2'$ block is
\[
 \begin{pmatrix}
 1+a_0^2\mu_0\nu&-a_0\mu_0\sqrt\nu\\
 -a_0\mu_0\sqrt\nu&\mu_0
 \end{pmatrix},\qquad
 \nu=(N^{-1})_0\ge0,\qquad \det=\mu_0>0.          \tag{R3E5}
\]
This includes $a$ infinitesimal, $r$ infinitesimal, and their
independent simultaneous opening. The $R_1$ pivot is $A_0$ because
$Cl^2\sim0$ and $h\sim1$; it does not require a new soft root.

If $g$ is infinitesimal, use the unscaled $R_2=R_2'/a$, as $a$
is a unit. Apply the exact R5.2--R5.4 centered-$m$ identity, keeping
\[
 (m-K)R_{mc}^2,\qquad
 R_{mc}=-\xi+uv\eta+(1-g)R_2+huT,
\]
for a sufficiently large positive real $K$. In the remainder
replace only the formal matrix coefficient $m$ by $K$ and add
$\tau H_m$, where $\tau=(m-K)g$ is bounded and $H_m$ is the exact
R5.4 matrix. The remaining source coefficients, especially
$\alpha_*$, stay at their actual values. This is a whole-matrix
identity, not a physical-source alteration. The same coefficient
bounds hold; all normal entries divided by $\sqrt C$ vanish.
The three separate residue pivots remain $1,A_0,\alpha_{*,0}$;
the last block becomes
\[
 \begin{pmatrix}
 1+K\nu&-(K+\tau_0)\sqrt\nu\\
 -(K+\tau_0)\sqrt\nu&K+2\tau_0
 \end{pmatrix},\qquad
 \det=K+2\tau_0-\tau_0^2\nu.                       \tag{R3E6}
\]
Since $0\le\nu\le2$, any real $K>2\tau_0^2+1$ makes it strict.
Thus WA10 applies to the entire local parent in both source cases,
for all independent infinitesimal or zero $u,v$, either sign and
all states. Restore the exact retained square in the $g$-small case.

\paragraph{Every finite remote testing chart.}
If both finite testing residues are nonzero, use the inherited
arbitrary-coefficient-scale noncommuting-residue theorem. It remains
to treat one nonzero finite angle residue. In these charts select
positive real lower bounds for original positive event coefficients
when indicated and retain their exact positive excess packets;
no stationary-contact assertion is made for the selected remainder.

Use the actual original $C$ event with
\[
       J_C=\sqrt C\,(uW-Y),\qquad d_C=C^{-1/2}\in\Inf.
                                                               \tag{R3E7}
\]
If $u$ is a bounded unit and $v$ infinitesimal or zero, clip $n,A,D$
at positive real lower bounds and keep their positive excesses and
the whole original $m,\alpha$ packets. Solve
$W=(Y+d_CJ_C)/u$. The selected complete remainder has bounded
coefficients in $(J_C,X,Y,Z)$ and residue
\[
 J_C^2+n_*(1+U)X^2+A_*(1+U)Y^2
                         +\tfrac12(X+uZ)^2\succ0,               \tag{R3E8}
\]
where $u$ on the right is its nonzero real residue and $U=u^2$.
The $D$ event has zero residue and causes no difficulty.

If $u$ is infinitesimal or zero, $v$ a bounded unit and $a$ a unit,
then $m$ has a positive real lower bound (also when $g$ is small).
Clip $n,m,A,D$ at positive bounds, retain all excesses and $\alpha$,
and solve $Y=uW-d_CJ_C$. The complete selected residue is
\[
 (1+V)J_C^2+(n_*+1/2)(1+V)X^2
       +m_*(vW-X)^2+D_*VZ^2\succ0,\quad V=v_0^2>0.              \tag{R3E9}
\]
This retains every state direction and the whole signed core.

Finally suppose $a$ is infinitesimal in the latter testing chart.
Then $g\sim1$, $\mu$ is a positive unit, and R3E1 gives
$A_0=t_0/2$, $D_0=1/(2t_0)$. Keep the additional legal roots
$S=aW$, $R=uW$ and use $Y=R-d_CJ_C$. Clip only $n$, retain its
excess and $\alpha$, and retain the actual bounded $A,D,\mu$.
The $m$ event is exactly
\[
 m(vW-X+uvY-uZ)^2
       =\mu(vS-aX+auvY-auZ)^2.                    \tag{R3E10}
\]
The entire bounded matrix on $(J_C,S,X,R,Z)$ has separate
positive pivots $1+V$, $\mu_0V$ and $(n_*+1/2)(1+V)$.
Its remaining $R,Z$ block is
\[
 \begin{pmatrix}
 [t_0(1+V)+V]/2&-v_0/2\\
 -v_0/2&V/(2t_0)
 \end{pmatrix},\qquad
 \det={V^2(t_0+1)\over4t_0}>0.                    \tag{R3E11}
\]
The dependent roots $R,S$ are allowed in a positive Gram; no inverse
of $u$ or $a$ is introduced. WA10 at the real testing residues
establishes every finite remote chart above.

\paragraph{Projective join and exact scope.}
The source estimates give positive real lower bounds for $n,A,C,D$.
Therefore the complete selected-event endpoint polynomials R4E14,
R4E15 and R4E16 apply verbatim: clip these four original event
coefficients, retain all their excesses and the entire $m,\alpha$
packets, and use the original degree-$(2,2)$ homogeneous parent.
Those exact identities cover respectively one infinite Alice angle,
one infinite Bob angle, and both infinite angles, including the
zero values in the other projective chart. Their strict Grams do
not depend on any R4-specific source relation. Thus every finite
and infinite testing pair is covered. Polynomial axes, both signs
and the real/imaginary state split give literal zeros and complex
states. All inverses and radicals used are pure-source or pure-test.

This admits exactly R3 with $1-r$ a positive bounded unit. It neither
claims the remaining $1-r$-infinitesimal sources empty nor changes
the historical weak-source partition. No numerical table is a
premise. The focused checker verifies the new source residues, full
local residue blocks and complete finite remote residues; unchanged
WA, centered-$m$, and homogeneous endpoint identities are imported.

\paragraph{Why the existing discrete symmetry does not remove R1--R3.}
In the proved cyclic order $(\lambda,D,C,A,\alpha,m,\beta,n)$,
the stabilizer of the omitted pair $\{\lambda,\beta\}=\{0,6\}$
inside the CHSH dihedral action consists of the identity and
$i\mapsto6-i$. The latter is exactly
\[
 (u,v;W,X,Y,Z;n,m,C,A,D,\alpha)
 \longmapsto(1/v,1/u;Z,X,Y,W;n,D,\alpha,A,m,C).
\]
It exchanges strong $C$/weak $\alpha$ with the opposite directed
orientation. Returning to the chosen weak orientation restores the
original source, so this known discrete action supplies no additional
within-orientation routing of R1--R3. This statement concerns only
the established dihedral action, not all conceivable analytical
changes of variables.
